% GENERATED FILE. Edit canonical lessons, then run npm run generate:books.
% canonical-source-hash: fnv1a64:b9150c06269b7a7d
% canonical-lessons: FR-C490-parrain, FR-C490-marraine, FR-C490-filleul, FR-C490-aine, FR-C490-cadet

\chapter{La famille au sens large}
\label{ch:fr-la-famille-au-sens-large}
\hlchaptermodality{490}

\begin{chapteropening}
\textbf{By the end of this chapter:} \emph{I can talk about godparents and older and younger siblings.}

Complete the last lesson of chapter 490: I can talk about godparents and older and younger siblings.
\end{chapteropening}

\section[le parrain]{le parrain — godfather; sponsor}

\label{lesson:FR-C490-parrain}



\begin{quote}
Before the new one: say the French for honeymoon, then the French for ceremony.
\end{quote}

\subsection*{You'll want to know: le parrain}
\textbf{le parrain} — \textquotedblleft{}godfather; sponsor\textquotedblright{}.

\begin{quote}
Mon parrain m'offre toujours un livre pour mon anniversaire.
\end{quote}

\begin{quote}
\emph{My godfather always gives me a book for my birthday.}
\end{quote}

Also someone who sponsors a newcomer: \emph{le parrain d'un projet}.

\begin{cousinweb}
\emph{le parrain} is inherited, through Old French \emph{parain}, from Vulgar Latin \emph{patrīnus} (EtymDB), built on \emph{pater} \textquotedblleft{}father\textquotedblright{} (Lewis \& Short).
\end{cousinweb}

\begin{grammarlens}[title={What you've built}]
One more word for telling, explaining and describing at B1.
\end{grammarlens}

\subsection*{Your turn}
\begin{itemize}
\raggedright
  \item \emph{Say it:} \emph{le parrain}
  \item \emph{Say it:} \emph{le parrain}, once more
  \item \emph{Say it:} a sentence of your own with \emph{le parrain}
  \item \emph{Recall:} say the French for honeymoon, then the French for ceremony, then say \emph{le parrain} again
\end{itemize}

\begin{tcolorbox}[breakable,colback=teal!4,colframe=teal!35!black,title={Before you move on}]
What does le parrain mean? (\textquotedblleft{}Godfather; sponsor\textquotedblright{}.) Say it once more.
\end{tcolorbox}
\section[la marraine]{la marraine — godmother}

\label{lesson:FR-C490-marraine}



\begin{quote}
Before the new one: say the French for ceremony, then the French for godfather.
\end{quote}

\subsection*{You'll want to know: la marraine}
\textbf{la marraine} — \textquotedblleft{}godmother\textquotedblright{}.

\begin{quote}
Ma marraine habite à la campagne.
\end{quote}

\begin{quote}
\emph{My godmother lives in the country.}
\end{quote}

Also a sponsor: \emph{la marraine d'une association}.

\begin{cousinweb}
\emph{la marraine} is inherited, through Old French \emph{marrene}, from Vulgar Latin \emph{mātrīna} (EtymDB), built on \emph{māter} \textquotedblleft{}mother\textquotedblright{} (Lewis \& Short).
\end{cousinweb}

\begin{grammarlens}[title={What you've built}]
One more word for telling, explaining and describing at B1.
\end{grammarlens}

\subsection*{Your turn}
\begin{itemize}
\raggedright
  \item \emph{Say it:} \emph{la marraine}
  \item \emph{Say it:} \emph{la marraine}, once more
  \item \emph{Say it:} a sentence of your own with \emph{la marraine}
  \item \emph{Recall:} say the French for ceremony, then the French for godfather, then say \emph{la marraine} again
\end{itemize}

\begin{tcolorbox}[breakable,colback=teal!4,colframe=teal!35!black,title={Before you move on}]
What does la marraine mean? (\textquotedblleft{}Godmother\textquotedblright{}.) Say it once more.
\end{tcolorbox}
\section[le filleul]{le filleul — godson, godchild}

\label{lesson:FR-C490-filleul}



\begin{quote}
Before the new one: say the French for godfather, then the French for godmother.
\end{quote}

\subsection*{You'll want to know: le filleul}
\textbf{le filleul} — \textquotedblleft{}godson, godchild\textquotedblright{}.

\begin{quote}
Elle emmène son filleul au zoo.
\end{quote}

\begin{quote}
\emph{She takes her godson to the zoo.}
\end{quote}

A goddaughter is \emph{la filleule}.

\begin{cousinweb}
\emph{le filleul} is inherited, through Old French \emph{fillol}, from Latin \emph{fīliolus} \textquotedblleft{}little son\textquotedblright{} (EtymDB), from \emph{fīlius} \textquotedblleft{}son\textquotedblright{} (Lewis \& Short).
\end{cousinweb}

\begin{grammarlens}[title={What you've built}]
One more word for telling, explaining and describing at B1.
\end{grammarlens}

\subsection*{Your turn}
\begin{itemize}
\raggedright
  \item \emph{Say it:} \emph{le filleul}
  \item \emph{Say it:} \emph{le filleul}, once more
  \item \emph{Say it:} a sentence of your own with \emph{le filleul}
  \item \emph{Recall:} say the French for godfather, then the French for godmother, then say \emph{le filleul} again
\end{itemize}

\begin{tcolorbox}[breakable,colback=teal!4,colframe=teal!35!black,title={Before you move on}]
What does le filleul mean? (\textquotedblleft{}Godson, godchild\textquotedblright{}.) Say it once more.
\end{tcolorbox}
\section[l'aîné]{l'aîné — eldest, older (sibling)}

\label{lesson:FR-C490-aine}



\begin{quote}
Before the new one: say the French for godmother, then the French for godson.
\end{quote}

\subsection*{You'll want to know: l'aîné}
\textbf{l'aîné} — \textquotedblleft{}eldest, older (sibling)\textquotedblright{}.

\begin{quote}
L'aîné de la famille a vingt ans.
\end{quote}

\begin{quote}
\emph{The eldest child in the family is twenty.}
\end{quote}

A woman is \emph{l'aînée}. \emph{Mon frère aîné} — my older brother.

\begin{cousinweb}
\emph{l'aîné} comes from Old French \emph{ainsné} (EtymDB), \textquotedblleft{}born before\textquotedblright{}: \emph{ains} \textquotedblleft{}before\textquotedblright{} goes back to Latin \emph{ante}, and \emph{né} to Latin \emph{nātus}, from \emph{nāscī} \textquotedblleft{}to be born\textquotedblright{} (Lewis \& Short).
\end{cousinweb}

\begin{grammarlens}[title={What you've built}]
One more word for telling, explaining and describing at B1.
\end{grammarlens}

\subsection*{Your turn}
\begin{itemize}
\raggedright
  \item \emph{Say it:} \emph{l'aîné}
  \item \emph{Say it:} \emph{l'aîné}, once more
  \item \emph{Say it:} a sentence of your own with \emph{l'aîné}
  \item \emph{Recall:} say the French for godmother, then the French for godson, then say \emph{l'aîné} again
\end{itemize}

\begin{tcolorbox}[breakable,colback=teal!4,colframe=teal!35!black,title={Before you move on}]
What does l'aîné mean? (\textquotedblleft{}Eldest, older (sibling)\textquotedblright{}.) Say it once more.
\end{tcolorbox}
\section[le cadet]{le cadet — younger (sibling); youngest}

\label{lesson:FR-C490-cadet}



\begin{quote}
Before the new one: say the French for godson, then the French for eldest.
\end{quote}

\subsection*{You'll want to know: le cadet}
\textbf{le cadet} — \textquotedblleft{}younger (sibling); youngest\textquotedblright{}.

\begin{quote}
Mon frère cadet a trois ans de moins que moi.
\end{quote}

\begin{quote}
\emph{My younger brother is three years younger than me.}
\end{quote}

A woman is \emph{la cadette}. \emph{C'est le cadet de mes soucis} — that is the least of my worries.

\begin{cousinweb}
\emph{le cadet} was borrowed from Gascon \emph{capdet} \textquotedblleft{}little chief\textquotedblright{}, a title for younger sons serving as officers; it goes back to Latin \emph{capitellum}, from \emph{caput} \textquotedblleft{}head\textquotedblright{} (Lewis \& Short).
\end{cousinweb}

\begin{grammarlens}[title={What you've built}]
One more word for telling, explaining and describing at B1.
\end{grammarlens}

\subsection*{Your turn}
\begin{itemize}
\raggedright
  \item \emph{Say it:} \emph{le cadet}
  \item \emph{Say it:} \emph{le cadet}, once more
  \item \emph{Say it:} a sentence of your own with \emph{le cadet}
  \item \emph{Recall:} say the French for godson, then the French for eldest, then say \emph{le cadet} again
\end{itemize}

\begin{tcolorbox}[breakable,colback=teal!4,colframe=teal!35!black,title={Before you move on}]
What does le cadet mean? (\textquotedblleft{}Younger (sibling); youngest\textquotedblright{}.) Say it once more.
\end{tcolorbox}
